% Luc Destombe --- licence CC BY-NC-SA 4.0
% https://creativecommons.org/licenses/by-nc-sa/4.0/deed.fr
% Fichier autonome : compiler avec pdflatex, deux fois (signets, nombre de pages).
\def\lyceeversion{1.0}
\documentclass[a4paper,10pt]{article}
\makeatletter
\newif\iflycee@niveaux
\lycee@niveauxtrue
\newif\iflycee@cm
\newif\iflycee@palatino

\RequirePackage[utf8]{inputenc}
\RequirePackage[T1]{fontenc}
\RequirePackage[french]{babel}
\RequirePackage{etoolbox}

\newcommand{\ShorthandsOff}{\@ifundefined{shorthandoff}{}{\shorthandoff{;:!?}}}
\newcommand{\ShorthandsOn}{\@ifundefined{shorthandon}{}{\shorthandon{;:!?}}}
\ShorthandsOff
\AtBeginDocument{\ShorthandsOn}
\AtBeginEnvironment{tikzpicture}{\ShorthandsOff}

\RequirePackage{geometry}
\RequirePackage{fancyhdr}
\RequirePackage{lastpage}
\RequirePackage{multicol}
\RequirePackage{array}
\RequirePackage{enumitem}
\RequirePackage{xcolor}
\RequirePackage{needspace}

\RequirePackage{amsmath,amssymb}
\iflycee@cm\else\RequirePackage{fourier}\fi
\DeclareMathAlphabet{\mathcal}{OMS}{cmsy}{m}{n}
\renewcommand{\ttdefault}{lmtt}
\begingroup\catcode`\_=\active \catcode`\^=\active
\global\let_\sb \global\let^\sp
\endgroup
\AtBeginDocument{\catcode`\_=12 \mathcode`\_="8000
  \catcode`\^=12 \mathcode`\^="8000 }
\newlength{\retraitcalculs}
\setlength{\retraitcalculs}{2em}

\RequirePackage{tikz}
\usetikzlibrary{arrows.meta,calc,babel,shapes.misc}
\RequirePackage[most]{tcolorbox}
\tcbuselibrary{skins,breakable}

\definecolor{coulDef}{HTML}{1F4E79}
\definecolor{coulProp}{HTML}{7030A0}
\definecolor{coulRem}{HTML}{C55A11}
\definecolor{coulEx}{HTML}{2E7D32}
\definecolor{coulExo}{HTML}{B03A2E}
\definecolor{coulPart}{HTML}{1F4E79}
\definecolor{coulMeth}{HTML}{1B6E4F}
\definecolor{coulCourbe}{HTML}{0B5ED7}
\definecolor{coulCourbeB}{HTML}{C00000}
\definecolor{coulCourbeC}{HTML}{2E7D32}
\definecolor{coulGrille}{HTML}{8C8C8C}
\definecolor{coulRep}{HTML}{1B6E4F}
\definecolor{coulBar}{HTML}{2E7D32}

\newcommand{\R}{\mathbb{R}}
\newcommand{\Rs}{\mathbb{R}^{*}}
\newcommand{\N}{\mathbb{N}}
\newcommand{\Ns}{\mathbb{N}^{*}}
\newcommand{\Z}{\mathbb{Z}}
\newcommand{\Q}{\mathbb{Q}}
\renewcommand{\le}{\leqslant}
\renewcommand{\ge}{\geqslant}
\renewcommand{\approx}{\simeq}
\newcommand{\vect}[1]{\overrightarrow{#1}}
\newcommand{\norme}[1]{\left\lVert #1 \right\rVert}
\newcommand{\intervalle}[2]{\left[#1\,;#2\right]}
\RequirePackage{cancel}
\renewcommand{\CancelColor}{\color{coulExo}}
\newcommand{\fois}[1]{\times{\color{coulExo}#1}}
\newcommand{\barre}[1]{\cancel{#1}}
\newlength{\lycee@eqh}\newlength{\lycee@eqd}
\newsavebox{\lycee@eqa}\newsavebox{\lycee@eqb}\newsavebox{\lycee@eqc}
\newcommand{\lycee@eqcol}[1]{\hbox{$\renewcommand{\arraystretch}{1.35}%
  \begin{array}[t]{@{}>{\displaystyle}r@{}>{\displaystyle{}}l@{}}#1\end{array}$}}
\newcommand{\lycee@eqmesure}[1]{%
  \ifdim\ht#1>\lycee@eqh\setlength{\lycee@eqh}{\ht#1}\fi
  \ifdim\dp#1>\lycee@eqd\setlength{\lycee@eqd}{\dp#1}\fi}
\newcommand{\lycee@eqtrait}{\hspace{1.2em}\tikz[baseline=0pt]\draw[dash pattern=on 2pt off 2pt,
  line width=0.5pt] (0,\lycee@eqh) -- (0,-\lycee@eqd);\hspace{1.2em}}
\newcommand{\eqproduit}[2]{\par\smallskip\noindent C'est une équation produit :
  $A\times B=0 \iff A=0$ ou $B=0$.\par\nopagebreak\smallskip
  \begingroup\setlength{\lycee@eqh}{0pt}\setlength{\lycee@eqd}{0pt}%
  \setbox\lycee@eqa=\lycee@eqcol{#1}\setbox\lycee@eqb=\lycee@eqcol{#2}%
  \lycee@eqmesure\lycee@eqa\lycee@eqmesure\lycee@eqb%
  \noindent\hspace*{\retraitcalculs}\usebox\lycee@eqa\lycee@eqtrait\usebox\lycee@eqb\endgroup\par\smallskip}
\newcommand{\eqproduittrois}[3]{\par\smallskip\noindent C'est une équation produit :
  $A\times B\times C=0 \iff A=0$ ou $B=0$ ou $C=0$.\par\nopagebreak\smallskip
  \begingroup\setlength{\lycee@eqh}{0pt}\setlength{\lycee@eqd}{0pt}%
  \setbox\lycee@eqa=\lycee@eqcol{#1}\setbox\lycee@eqb=\lycee@eqcol{#2}\setbox\lycee@eqc=\lycee@eqcol{#3}%
  \lycee@eqmesure\lycee@eqa\lycee@eqmesure\lycee@eqb\lycee@eqmesure\lycee@eqc%
  \noindent\hspace*{\retraitcalculs}\usebox\lycee@eqa\lycee@eqtrait\usebox\lycee@eqb\lycee@eqtrait\usebox\lycee@eqc\endgroup
  \par\smallskip}

\tikzset{grille/.style={coulGrille,line width=0.4pt}}
\newcommand{\quadrillage}[4]{\draw[grille,step=1] (#1,#2) grid (#3,#4);}
\tikzset{
  croix/.style={cross out,draw,line width=0.6pt,inner sep=0pt,minimum size=4pt},
  croixdroite/.style={croix,rotate=45,line width=0.8pt},
}
\newcommand{\pt}[3]{\path (#1) node[croix]{} node[#2,font=\small,inner sep=2.5pt]{$#3$};}

\newcommand{\lycee@repere}[4]{%
  \draw[grille,step=1] (#1,#2) grid (#3,#4);
  \draw[-{Stealth[length=2mm]},thick] (#1,0)--({#3+0.4},0) node[below left=-1pt]{$x$};
  \draw[-{Stealth[length=2mm]},thick] (0,#2)--(0,{#4+0.4}) node[below left=-1pt]{$y$};
  \node[below left,font=\scriptsize,inner sep=1.5pt] at (0,0) {$0$};}
\newcommand{\lycee@gradx}[1]{\foreach \k/\l in {#1}{%
  \draw (\k,0.07)--(\k,-0.07) node[below,font=\scriptsize,inner sep=1.5pt]{$\l$};}}
\newcommand{\lycee@grady}[1]{\foreach \k/\l in {#1}{%
  \draw (0.07,\k)--(-0.07,\k) node[left,font=\scriptsize,inner sep=1.5pt]{$\l$};}}
\newcommand{\lycee@intff}[2]{\left[#1\,;#2\right]}
\newcommand{\lycee@intfo}[2]{\left[#1\,;#2\right[}
\newcommand{\lycee@intof}[2]{\left]#1\,;#2\right]}
\newcommand{\lycee@intoo}[2]{\left]#1\,;#2\right[}
\newcommand{\lycee@ens}[1]{\left\{#1\right\}}
\AtBeginDocument{%
  \providecommand{\repere}{\lycee@repere}%
  \providecommand{\gradx}{\lycee@gradx}%
  \providecommand{\grady}{\lycee@grady}%
  \providecommand{\Cf}{\mathcal{C}\sb{f}}%
  \providecommand{\Cg}{\mathcal{C}\sb{g}}%
  \providecommand{\intff}{\lycee@intff}%
  \providecommand{\intfo}{\lycee@intfo}%
  \providecommand{\intof}{\lycee@intof}%
  \providecommand{\intoo}{\lycee@intoo}%
  \providecommand{\ens}{\lycee@ens}}

\newlist{questions}{enumerate}{3}
\setlist[questions,1]{label=\arabic*\textdegree),leftmargin=*,itemsep=2pt,topsep=3pt}
\setlist[questions,2]{label=\alph*),leftmargin=*,itemsep=1pt,topsep=2pt}
\setlist[questions,3]{label=\roman*),leftmargin=*,itemsep=1pt,topsep=1pt}
\setlist[itemize]{label=\textbullet,leftmargin=*,itemsep=2pt,topsep=3pt}

\newcommand{\besoinplace}[1]{\par\penalty\@M\begingroup
  \dimen@=#1\relax\dimen@ii=\pagegoal\advance\dimen@ii-\pagetotal
  \ifdim\dimen@>\dimen@ii\ifdim\dimen@ii>\z@\vfil\fi\break\fi\endgroup}

\newcommand{\lycee@pied}{\thepage/\pageref{LastPage}}
\newcommand{\entete}[2]{%
  \pagestyle{fancy}\fancyhf{}%
  \fancyhead[L]{\footnotesize\color{coulPart}\textbf{#1}}%
  \fancyhead[R]{\footnotesize\color{coulPart}\textbf{#2}}%
  \fancyfoot[C]{\footnotesize\lycee@pied}%
  \renewcommand{\headrulewidth}{0.4pt}%
  \renewcommand{\headrule}{\hbox to\headwidth{%
    \color{coulPart!40}\leaders\hrule height \headrulewidth\hfill}}}

\RequirePackage{ccicons}
\newcommand{\lycee@licence}{{\scriptsize\color{gray}%
  \copyright\ Luc Destombe --- \ccbyncsa\ CC BY-NC-SA 4.0}}
\newif\iflycee@licence \lycee@licencetrue
\newcommand{\sanslicence}{\lycee@licencefalse}
\AtBeginDocument{\iflycee@licence\AddToHookNext{shipout/foreground}{%
  \put(0,0){\rlap{\kern\dimexpr1in+\hoffset+\oddsidemargin+\textwidth\relax
    \llap{\raisebox{-\dimexpr1in+\voffset+\topmargin+\headheight+\headsep
      +\textheight+\footskip\relax}{\lycee@licence}}}}}\fi}

\newcommand{\lycee@site}{faire-des-maths.fr}
\newcommand{\lycee@version}{{\scriptsize\color{gray}\lycee@site
  \ifdefined\lyceeversion\ --- Version~\lyceeversion\fi}}
\AtBeginDocument{\iflycee@licence\AddToHookNext{shipout/foreground}{%
  \put(0,0){\rlap{\kern\dimexpr1in+\hoffset+\oddsidemargin\relax
    \raisebox{-\dimexpr1in+\voffset+\topmargin+\headheight+\headsep
      +\textheight+\footskip\relax}{\lycee@version}}}}\fi}

\newcommand{\formulesserrees}{%
  \setlength{\abovedisplayskip}{3pt}\setlength{\belowdisplayskip}{3pt}%
  \setlength{\abovedisplayshortskip}{2pt}\setlength{\belowdisplayshortskip}{3pt}}

\setlength{\parindent}{0pt}

\RequirePackage[implicit=false,hidelinks,pdfencoding=auto,psdextra,bookmarksopen,
  bookmarksopenlevel=1,pdfcreator={},pdfproducer={}]{hyperref}
\RequirePackage{bookmark}
\hypersetup{pdfauthor={Luc Destombe},pdfkeywords={CC BY-NC-SA 4.0}}
\newcounter{lycee@signet}
\newcommand{\signet}[3][]{\stepcounter{lycee@signet}%
  \edef\lycee@dest{\if\relax\detokenize{#1}\relax
    signet.\arabic{lycee@signet}\else#1\fi}%
  \hypertarget{\lycee@dest}{}\bookmark[level=#2,dest=\lycee@dest]{#3}}
\pdfstringdefDisableCommands{%
  \def\R{R}\def\N{N}\def\Z{Z}\def\Q{Q}\def\vect#1{#1}\def\Cf{Cf}\def\Cg{Cg}\def\fois#1{×#1}%
  \def\barre#1{#1}\def\intervalle#1#2{[#1;#2]}\def\intff#1#2{[#1;#2]}%
  \def\textsuperscript#1{#1}\def\dfrac#1#2{#1/#2}\def\frac#1#2{#1/#2}%
  \def\vec#1{#1⃗}}%
\RequirePackage[official]{eurosym}

\geometry{top=0.8cm,bottom=1.3cm,left=1cm,right=1cm,footskip=0.6cm}

\pagestyle{fancy}
\fancyhf{}
\renewcommand{\headrulewidth}{0pt}
\fancyfoot[C]{\footnotesize Page \thepage{} / \pageref{LastPage}}

\newcommand{\titrefiche}[2]{%
  \begin{center}
    \Large\textbf{#1}\\[2pt]
    \large\textbf{#2}\\[2pt]
  \end{center}
  \vspace{0.10cm}}

\setlength{\columnseprule}{0.4pt}
\setlength{\columnsep}{15pt}
\renewcommand{\columnseprulecolor}{\color{black!45}}
\setlength{\multicolsep}{2pt}

\newcounter{partiefiche}
\renewcommand{\thepartiefiche}{\Roman{partiefiche}}
\newif\ifapresbandeau
\newcommand{\lycee@bandeau}[1]{%
  \par\penalty-600
  \vspace{0.08cm}%
  \noindent\colorbox{coulPart}{%
    \parbox{\dimexpr\linewidth-2\fboxsep\relax}{%
      \color{white}\bfseries\raggedright\strut#1\strut}}%
  \nopagebreak\par\nopagebreak
  \vspace{0.06cm}\nopagebreak
  \apresbandeautrue}
\newcommand{\partiefiche}[1]{%
  \refstepcounter{partiefiche}%
  \lycee@bandeau{\signet{1}{\thepartiefiche{} --- #1}\thepartiefiche{} --- #1}}
\newcommand{\partiefichesansnum}[1]{\lycee@bandeau{\signet{1}{#1}#1}}

\newcounter{exo}
\newlength{\espaceexo}\setlength{\espaceexo}{7pt}
\newlength{\espacetitre}\setlength{\espacetitre}{2pt}
\newcommand{\exo}[1][\relax]{%
  \par
  \ifapresbandeau\apresbandeaufalse\else\penalty-150\addvspace{\espaceexo}\fi
  \refstepcounter{exo}%
  \noindent\signet[exercice-\arabic{exo}]{2}{Exercice \arabic{exo}}%
  \textbf{\color{coulExo}\underline{Exercice \theexo}}%
  \ifx\relax#1\relax\else\textbf{ : #1}\fi
  \par\nopagebreak\vspace{\espacetitre}\nopagebreak
  \ignorespaces}

\setlist[enumerate]{nosep,leftmargin=*,itemsep=2pt,topsep=2pt}
\setlist[itemize]{nosep,leftmargin=*,topsep=2pt}
\newcommand{\choix}[3]{\par\hspace*{1em}%
  \makebox[0.3\linewidth][l]{a) #1}\makebox[0.3\linewidth][l]{b) #2}c) #3}
\newcommand{\deux}[2]{\par\hspace*{0.6em}%
  \makebox[0.48\linewidth][l]{#1}#2}
\newcommand{\trois}[3]{\par\hspace*{0.6em}%
  \makebox[0.32\linewidth][l]{#1}\makebox[0.32\linewidth][l]{#2}#3}

  \renewcommand{\exo}[1][\relax]{%
    \ifapresbandeau\apresbandeaufalse\else\par\penalty-150\fi
    \refstepcounter{exo}%
    \vspace{0.05cm}%
    \noindent\signet[exercice-\arabic{exo}]{2}{Exercice \arabic{exo}}%
    \textbf{\color{coulExo}\underline{Exercice \theexo}}%
    \ifx\relax#1\relax\else\textbf{ : #1}\fi%
    \par\nopagebreak
    \ignorespaces}
  \newcommand{\rappel}[1]{%
    \par\vspace{0.06cm}\nopagebreak
    \noindent\fcolorbox{black!35}{black!5}{%
      \parbox{\dimexpr\linewidth-2\fboxsep-2\fboxrule\relax}{%
        \small\textbf{Rappel.} #1}}%
    \par\vspace{0.06cm}\nopagebreak}
  \setlist[enumerate]{nosep,leftmargin=*}
  \setlist[itemize]{nosep,leftmargin=*}
  \setlength{\multicolsep}{3pt plus 1pt minus 1pt}
  \setlength{\premulticols}{10pt}
  \setlength{\postmulticols}{10pt}
  \newlist{expr}{enumerate}{1}
  \setlist[expr]{label={\Alph*\ $=$},nosep,leftmargin=*,itemsep=0pt}

\renewenvironment{center}{\par\addvspace{3pt}\centering}{\par\addvspace{3pt}}
\formulesserrees
\AtBeginDocument{\formulesserrees}
\clubpenalty=10000
\widowpenalty=10000
\displaywidowpenalty=10000
\brokenpenalty=10000
\setlength{\parskip}{0pt}

\RequirePackage{ragged2e}
\setlength{\RaggedRightRightskip}{0pt plus 1fil}
\AtBeginDocument{\RaggedRight\binoppenalty=10000 \relpenalty=10000 }
\g@addto@macro\@parboxrestore{\RaggedRight}
\makeatother
\geometry{top=0.8cm,bottom=1.3cm,left=1cm,right=1cm,footskip=0.7cm,headheight=12pt}

\newtcblisting{python}{%
  listing only,enhanced,breakable,
  colback=black!3,colframe=black!35,boxrule=0.5pt,arc=1pt,
  left=6pt,right=6pt,top=2pt,bottom=2pt,
  listing options={language=Python,basicstyle=\ttfamily\small,
    keywordstyle=\color{coulMeth}\bfseries,columns=fullflexible,
    keepspaces=true,showstringspaces=false,aboveskip=0pt,belowskip=0pt}}

\begin{document}
\titrefiche{Terminale spécialité --- Suites et récurrence}{Fiche d'exercices du chapitre}

\begin{multicols}{2}

\partiefiche{Suites arithmétiques et géométriques (rappels de 1\textsuperscript{re})}

\exo[Reconnaître une suite arithmétique]
Pour chacune des suites $(a_{n})$ à $(d_{n})$ ci-dessous, dire si elle est arithmétique ;
si oui, donner son premier terme et sa raison.
\begin{enumerate}[label=\alph*)]
  \item $(a_{n})$ définie pour tout $n\in\N$ par $a_{n}=5n-4$.
  \item $(b_{n})$ définie pour tout $n\in\N$ par $b_{n}=n^{2}+2$.
  \item $(c_{n})$ définie par $c_{0}=3$ et, pour tout $n\in\N$, $c_{n+1}=2c_{n}-1$.
  \item $(d_{n})$ définie pour tout $n\in\N$ par $d_{n}=(n+3)^{2}-n^{2}$.
\end{enumerate}

\exo[Reconnaître une suite géométrique]
Pour chacune des suites $(a_{n})$ à $(d_{n})$ ci-dessous, dire si elle est géométrique ;
si oui, donner son premier terme et sa raison.
\begin{enumerate}[label=\alph*)]
  \item $(a_{n})$ définie pour tout $n\in\N$ par $a_{n}=6^{n}$.
  \item $(b_{n})$ définie pour tout $n\in\N$ par $b_{n}=n^{3}$.
  \item $(c_{n})$ définie pour tout $n\in\N$ par $c_{n}=-3\times 2^{n+2}$.
  \item $(d_{n})$ définie par $d_{0}=2$ et, pour tout $n\in\N$, $d_{n+1}=2d_{n}+4n$.
\end{enumerate}

\exo[Deux modèles pour une durée d'écoute]
En janvier 2026, Inès a écouté $30$~h de musique en ligne, et Noé $16$~h. On suppose
que, chaque mois, la durée d'écoute d'Inès augmente de $1{,}5$~h et celle de Noé de
$8\,\%$. Pour tout $n\in\Ns$, on note $a_{n}$ et $b_{n}$ les durées d'écoute, en heure,
d'Inès et de Noé le $n$-ième mois : $a_{1}=30$ et $b_{1}=16$.
\begin{enumerate}
  \item Calculer $a_{2}$, $a_{3}$, $b_{2}$ et $b_{3}$.
  \item Exprimer $a_{n+1}$ en fonction de $a_{n}$, et $b_{n+1}$ en fonction de $b_{n}$ ;
        en déduire la nature de chaque suite.
  \item Exprimer $a_{n}$ et $b_{n}$ en fonction de $n$.
  \item Selon ce modèle, combien de temps chacun écoutera-t-il de la musique en
        décembre 2026 ?
  \item Compléter la fonction Python ci-dessous pour qu'elle renvoie $a_{12}$ et
        $b_{12}$.
\begin{python}
def ecoute():
    a = ...
    b = ...
    for i in range(...):
        a = ...
        b = ...
    return a, b
\end{python}
  \item À l'aide de la calculatrice, déterminer le premier mois où Noé écoutera plus de
        musique qu'Inès.
\end{enumerate}

\exo[Vrai ou faux : le contre-exemple]
Pour chaque affirmation, dire si elle est vraie ou fausse, en justifiant. Pour montrer
qu'une affirmation est fausse, un contre-exemple suffit.
\begin{enumerate}[label=\alph*)]
  \item La suite $(u_{n})$ définie pour tout $n\in\N$ par $u_{n}=n^{2}+1$ est
        arithmétique.
  \item Si une suite $(v_{n})$ vérifie $v_{1}-v_{0}=v_{2}-v_{1}$, alors elle est
        arithmétique.
  \item La somme de deux suites géométriques est une suite géométrique.
  \item Le produit de deux suites géométriques est une suite géométrique.
\end{enumerate}

\exo[Une plateforme qui perd des abonnés]
Une plateforme de vidéos compte $2{,}4$ millions d'abonnés en janvier 2026. Chaque
année, elle en perd $6\,\%$. Pour tout $n\in\N$, on note $u_{n}$ le nombre d'abonnés,
en million, en janvier $2026+n$ ; ainsi $u_{0}=2{,}4$.
\begin{enumerate}
  \item Calculer $u_{1}$.
  \item Exprimer $u_{n+1}$ en fonction de $u_{n}$ ; en déduire la nature de la suite
        $(u_{n})$, puis $u_{n}$ en fonction de $n$.
  \item Combien la plateforme comptera-t-elle d'abonnés en janvier 2030 ?
  \item On considère la fonction Python ci-dessous. Que représente la valeur qu'elle
        renvoie ?
\begin{python}
def seuil():
    n = 0
    u = 2.4
    while u >= 1.2:
        n = n + 1
        u = 0.94 * u
    return n
\end{python}
  \item À l'aide de la calculatrice, déterminer cette valeur, puis l'année
        correspondante.
\end{enumerate}

\partiefiche{Suites auxiliaires}

\exo[Calcul : écrire sous la forme d'un seul quotient]
Pour chaque expression, $x$ est un nombre réel pour lequel elle est définie. L'écrire
sous la forme d'un seul quotient.
\begin{center}
$A=2+\dfrac{3}{x+1}$ \qquad\qquad $B=4-\dfrac{6}{2-x}$ \\[6pt]
$C=\dfrac{2}{3}+\dfrac{1}{2x+1}$ \qquad\qquad $D=3+\dfrac{x+2}{x-1}$ \\[6pt]
$E=\dfrac{3x}{x+2}-\dfrac{4}{x+2}$ \qquad\qquad $F=\dfrac{2}{x-1}-\dfrac{3}{x+2}$
\end{center}

\exo[Calcul : exprimer $u$ en fonction de $v$]
Dans chaque cas, $u$ et $v$ sont deux nombres réels pour lesquels les quotients sont
définis. Exprimer $u$ en fonction de $v$.
\begin{center}
a)~$v=\dfrac{u}{u-1}$ \qquad\qquad b)~$v=\dfrac{u+1}{u-3}$ \\[6pt]
c)~$v=\dfrac{u-2}{u+1}$ \qquad\qquad d)~$v=\dfrac{1}{u+3}$
\end{center}

\exo[Les abonnés d'une médiathèque]
En septembre 2026, une médiathèque compte $4\,000$ abonnés. Chaque année, $15\,\%$ des
abonnés ne renouvellent pas leur abonnement, et $450$ nouveaux abonnés s'inscrivent.
Pour tout $n\in\N$, on note $u_{n}$ le nombre d'abonnés en septembre $2026+n$ ; ainsi
$u_{0}=4\,000$.
\begin{enumerate}
  \item Calculer $u_{1}$ et $u_{2}$.
  \item Justifier que, pour tout $n\in\N$, on a $u_{n+1}=0{,}85\,u_{n}+450$.
  \item Compléter la fonction Python ci-dessous pour qu'elle renvoie $u_{n}$.
\begin{python}
def terme(n):
    u = ...
    for i in range(...):
        u = ...
    return u
\end{python}
  \item Soit la suite $(v_{n})$ définie pour tout $n\in\N$ par $v_{n}=u_{n}-3\,000$.
        Démontrer que $(v_{n})$ est géométrique ; préciser sa raison et son premier
        terme.
  \item Exprimer $v_{n}$, puis $u_{n}$, en fonction de $n$.
  \item Justifier que le nombre d'abonnés restera toujours supérieur à $3\,000$.
\end{enumerate}

\exo[Une suite auxiliaire qui dépend de $n$]
Soit la suite $(u_{n})$ définie par $u_{0}=1$ et, pour tout $n\in\N$,
$u_{n+1}=3u_{n}-4n+6$.
\begin{enumerate}
  \item Calculer $u_{1}$ et $u_{2}$.
  \item Écrire une fonction Python \texttt{suite(n)} qui renvoie $u_{n}$.
  \item Soit la suite $(v_{n})$ définie pour tout $n\in\N$ par $v_{n}=u_{n}-2n+2$.
        Démontrer que $(v_{n})$ est géométrique de raison $3$.
  \item En déduire l'expression de $u_{n}$ en fonction de $n$.
\end{enumerate}

\exo[Une suite auxiliaire géométrique définie par un quotient]
Soit la suite $(u_{n})$ définie par $u_{0}=0$ et, pour tout $n\in\N$,
$u_{n+1}=\dfrac{3u_{n}+2}{u_{n}+2}$. On admet que $u_{n}\geq 0$ pour tout $n\in\N$.
Soit la suite $(v_{n})$ définie pour tout $n\in\N$ par $v_{n}=\dfrac{u_{n}-2}{u_{n}+1}$.
\begin{enumerate}
  \item Calculer $u_{1}$ et $u_{2}$, puis $v_{0}$, $v_{1}$ et $v_{2}$.
  \item Démontrer que la suite $(v_{n})$ est géométrique ; préciser sa raison.
  \item Exprimer $v_{n}$ en fonction de $n$.
  \item Exprimer $u_{n}$ en fonction de $v_{n}$, puis en fonction de $n$.
\end{enumerate}

\exo[Une suite auxiliaire arithmétique définie par un quotient]
Soit la suite $(u_{n})$ définie par $u_{0}=1$ et, pour tout $n\in\N$,
$u_{n+1}=\dfrac{u_{n}-1}{u_{n}+3}$. On admet que $u_{n}>-1$ pour tout $n\in\N$.
Soit la suite $(v_{n})$ définie pour tout $n\in\N$ par $v_{n}=\dfrac{1}{u_{n}+1}$.
\begin{enumerate}
  \item Calculer $u_{1}$ et $u_{2}$, puis $v_{0}$, $v_{1}$ et $v_{2}$.
  \item Démontrer que la suite $(v_{n})$ est arithmétique de raison $\frac{1}{2}$.
  \item En déduire l'expression de $v_{n}$, puis de $u_{n}$, en fonction de $n$.
  \item Démontrer que, pour tout $n\in\N$, on a
        $u_{n+1}-u_{n}=-\dfrac{(u_{n}+1)^{2}}{u_{n}+3}$. En déduire que
        $u_{n+1}\leq u_{n}$.
\end{enumerate}

\partiefiche{Suites définies par une fonction}

\exo[Représenter une suite et conjecturer]
Soit $f$ la fonction définie sur $\R$ par $f(x)=0{,}5x+3$ et la suite $(u_{n})$ définie
par $u_{0}=0$ et, pour tout $n\in\N$, $u_{n+1}=f(u_{n})$.
\begin{enumerate}
  \item Calculer $u_{1}$, $u_{2}$ et $u_{3}$.
  \item Donner la relation de récurrence de $(u_{n})$.
  \item Dans un repère orthonormé, tracer la droite d'équation $y=x$ et la courbe de $f$,
        puis placer $u_{0}$, $u_{1}$, $u_{2}$ et $u_{3}$ sur l'axe des abscisses.
  \item Conjecturer le sens de variation de $(u_{n})$ et le nombre dont $u_{n}$ se
        rapproche quand $n$ devient grand.
  \item Reprendre les questions 3 et 4 avec la suite $(v_{n})$ définie par $v_{0}=10$
        et, pour tout $n\in\N$, $v_{n+1}=f(v_{n})$.
\end{enumerate}

\exo[Explicite ou par récurrence]
Soit $f$ la fonction définie sur $\R$ par $f(x)=x^{2}-4x+5$. Soit la suite $(u_{n})$
définie pour tout $n\in\N$ par $u_{n}=f(n)$, et la suite $(v_{n})$ définie par
$v_{0}=1$ et, pour tout $n\in\N$, $v_{n+1}=f(v_{n})$.
\begin{enumerate}
  \item Calculer $u_{0}$, $u_{1}$, $u_{2}$, $u_{3}$, puis $u_{10}$.
  \item Calculer $v_{1}$, $v_{2}$, $v_{3}$ et $v_{4}$.
  \item Peut-on calculer $v_{10}$ en remplaçant $n$ par $10$ dans une formule ?
        Déterminer $v_{10}$.
  \item Soit la suite $(w_{n})$ définie par $w_{0}=3$ et, pour tout $n\in\N$,
        $w_{n+1}=f(w_{n})$. Calculer $w_{1}$, $w_{2}$ et $w_{3}$ : que remarque-t-on ?
\end{enumerate}

\exo[Lire une suite sur un graphique]
Soit $f$ la fonction définie sur $\intfo{0}{+\infty}$ par $f(x)=\sqrt{2x+3}$. Sa courbe
$\Cf$ et la droite d'équation $y=x$ sont tracées ci-dessous.
\begin{center}
\begin{tikzpicture}[x=0.75cm,y=0.75cm]
  \repere{0}{0}{8}{5}
  \gradx{1,2,3,4,5,6,7,8} \grady{1,2,3,4,5}
  \draw[coulCourbeB,line width=0.8pt] (0,0)--(5,5) node[right,font=\footnotesize]{$y=x$};
  \draw[coulCourbe,line width=1.1pt] plot[domain=0:8,samples=80] (\x,{sqrt(2*\x+3)});
  \node[coulCourbe,font=\footnotesize] at (7.4,4.75) {$\Cf$};
\end{tikzpicture}
\end{center}
Soit les suites $(u_{n})$ et $(v_{n})$ définies par $u_{0}=0$, $v_{0}=6$ et, pour tout
$n\in\N$, $u_{n+1}=f(u_{n})$ et $v_{n+1}=f(v_{n})$.
\begin{enumerate}
  \item Calculer $u_{1}$ et $v_{1}$ (valeurs exactes, puis arrondies au centième).
  \item Sur le graphique, placer $u_{0}$, $u_{1}$, $u_{2}$, $u_{3}$, puis $v_{0}$,
        $v_{1}$, $v_{2}$, $v_{3}$ sur l'axe des abscisses.
  \item Conjecturer le sens de variation de chaque suite et le nombre dont ses termes se
        rapprochent.
  \item Résoudre dans $\intfo{0}{+\infty}$ l'équation $f(x)=x$. Ce résultat
        est-il cohérent avec la conjecture ?
\end{enumerate}

\partiefiche{Sens de variation : trois méthodes}

\exo[Étudier le sens de variation]
Étudier le sens de variation de chacune des suites $(a_{n})$ à $(e_{n})$ ci-dessous.
\begin{enumerate}[label=\alph*)]
  \item $(a_{n})$ définie pour tout $n\in\N$ par $a_{n}=2n^{2}-9n+1$.
  \item $(b_{n})$ définie pour tout $n\in\Ns$ par $b_{n}=\dfrac{2^{n}}{n}$.
  \item $(c_{n})$ définie pour tout $n\in\N$ par $c_{n}=n^{4}+3n^{2}$.
  \item $(d_{n})$ définie pour tout $n\in\N$ par $d_{n}=7-\dfrac{4}{n+1}$.
  \item $(e_{n})$ définie pour tout $n\in\N$ par $e_{n}=5\times 0{,}6^{n}$.
\end{enumerate}

\exo[Une suite définie par un quotient]
Soit la suite $(u_{n})$ définie pour tout $n\in\N$ par $u_{n}=\dfrac{2n+1}{n+3}$.
\begin{enumerate}
  \item Calculer $u_{0}$, $u_{1}$ et $u_{2}$.
  \item Démontrer que, pour tout $n\in\N$, on a
        $u_{n+1}-u_{n}=\dfrac{5}{(n+3)(n+4)}$.
  \item En déduire le sens de variation de $(u_{n})$.
  \item Vérifier que, pour tout $n\in\N$, on a $u_{n}=2-\dfrac{5}{n+3}$. En déduire que
        $u_{n}<2$.
  \item Retrouver le sens de variation de $(u_{n})$ avec la fonction $g$ définie sur
        $\intfo{0}{+\infty}$ par $g(x)=2-\dfrac{5}{x+3}$.
\end{enumerate}

\exo[Décroissante à partir d'un rang]
Soit la suite $(u_{n})$ définie pour tout $n\in\N$ par $u_{n}=n\times 0{,}9^{n}$.
\begin{enumerate}
  \item À l'aide de la calculatrice, donner les valeurs arrondies au millième de
        $u_{0}$ à $u_{12}$. Quel semble être le sens de variation de $(u_{n})$ ?
  \item Démontrer que, pour tout $n\in\Ns$, on a
        $\dfrac{u_{n+1}}{u_{n}}=\dfrac{0{,}9\,(n+1)}{n}$.
  \item Résoudre dans $\Ns$ l'inéquation $\dfrac{0{,}9\,(n+1)}{n}\leq 1$.
  \item En déduire le sens de variation de $(u_{n})$, puis la plus grande valeur prise
        par $u_{n}$.
\end{enumerate}

\exo[Suites récurrentes : la méthode de la différence]
Étudier le sens de variation de chacune des suites $(u_{n})$, $(v_{n})$, $(w_{n})$
ci-dessous.
\begin{enumerate}[label=\alph*)]
  \item $(u_{n})$ définie par $u_{0}=0$ et, pour tout $n\in\N$,
        $u_{n+1}=u_{n}+n^{2}-5n+6$ (on pourra factoriser $n^{2}-5n+6$).
  \item $(v_{n})$ définie par $v_{0}=1$ et, pour tout $n\in\N$,
        $v_{n+1}=v_{n}-\dfrac{2}{v_{n}^{2}+1}$.
  \item $(w_{n})$ définie par $w_{0}=2$ et, pour tout $n\in\N$,
        $w_{n+1}=w_{n}^{2}-w_{n}+1$.
\end{enumerate}

\partiefiche{Le raisonnement par récurrence}

\exo[Comprendre une propriété $P_{n}$]
Soit la suite $(u_{n})$ définie par $u_{0}=1$ et, pour tout $n\in\N$,
$u_{n+1}=3u_{n}+2$.
\begin{enumerate}
  \item Calculer $u_{1}$ et $u_{2}$.
  \item Compléter la fonction Python ci-dessous pour que \texttt{suite(n)} renvoie
        $u_{n}$.
\begin{python}
def suite(n):
    u = ...
    for i in range(...):
        u = ...
    return u
\end{python}
  \item Pour tout $n\in\N$, on note $P_{n}$ la propriété «~$u_{n}\geq 1$~», $Q_{n}$ la
        propriété «~$u_{n}=4n+1$~», $R_{n}$ la propriété «~$u_{n}=2\times 3^{n}-1$~»
        et $S_{n}$ la propriété «~$0<u_{n}<10$~». Pour chacune, dire si elle est vraie
        pour $n=0$, pour $n=1$ et pour $n=2$.
  \item Peut-on affirmer que, pour tout $n\in\N$, $Q_{n}$ est vraie ? Et $S_{n}$ ?
  \item Les trois premiers rangs suffisent-ils pour affirmer que $R_{n}$ est vraie pour
        tout $n\in\N$ ?
\end{enumerate}

\exo[L'initialisation est indispensable]
Pour tout $n\in\N$, on note $P_{n}$ la propriété «~$5^{n}+2$ est un multiple
de~$4$~».
\begin{enumerate}
  \item Soit $n\in\N$ tel que $P_{n}$ est vraie : il existe un entier $k$ tel que
        $5^{n}+2=4k$. Démontrer que $5^{n+1}+2=4(5k-2)$, puis que $P_{n+1}$ est vraie.
  \item La propriété $P_{0}$ est-elle vraie ?
  \item Peut-on affirmer que, pour tout $n\in\N$, $5^{n}+2$ est un multiple de $4$ ?
  \item Démontrer par récurrence que, pour tout $n\in\N$, $5^{n}-1$ est un multiple de
        $4$.
\end{enumerate}

\exo[Démontrer une inégalité par récurrence]
Pour chaque suite, démontrer l'inégalité par récurrence.
\begin{enumerate}[label=\alph*)]
  \item $(a_{n})$ définie par $a_{0}=4$ et, pour tout $n\in\N$, $a_{n+1}=2a_{n}-3$ :
        pour tout $n\in\N$, on a $a_{n}\geq 3$.
  \item $(b_{n})$ définie par $b_{0}=1$ et, pour tout $n\in\N$,
        $b_{n+1}=0{,}5\,b_{n}+1$ : pour tout $n\in\N$, on a $b_{n}\leq 2$.
  \item $(c_{n})$ définie par $c_{0}=0$ et, pour tout $n\in\N$, $c_{n+1}=3c_{n}+1$ :
        pour tout $n\in\N$, on a $c_{n}\geq 0$.
  \item $(d_{n})$ définie par $d_{0}=2$ et, pour tout $n\in\N$, $d_{n+1}=d_{n}^{2}+1$ :
        pour tout $n\in\N$, on a $d_{n}\geq 2$.
\end{enumerate}

\exo[Comparer une suite à $2n^{2}$]
Soit la suite $(u_{n})$ définie par $u_{0}=2$ et, pour tout $n\in\N$,
$u_{n+1}=u_{n}+4n+3$.
\begin{enumerate}
  \item Calculer $u_{1}$, $u_{2}$ et $u_{3}$, et les comparer à $2n^{2}$.
  \item Démontrer par récurrence que, pour tout $n\in\N$, on a $u_{n}\geq 2n^{2}$.
  \item En déduire que $u_{n}\geq 200$ pour tout entier $n\geq 10$.
\end{enumerate}

\exo[Une inégalité à partir d'un rang]
\begin{enumerate}
  \item Résoudre dans $\N$ l'inéquation $4n+2\geq 2n+3$.
  \item Pour $n=0$, $1$, $2$ et $3$, dire si l'inégalité $2^{n}\geq 2n+1$ est vraie.
  \item Démontrer par récurrence que, pour tout entier $n\geq 3$, on a
        $2^{n}\geq 2n+1$.
  \item L'inégalité est vraie pour $n=0$. Pourquoi ne peut-on pas faire
        l'initialisation au rang $0$ ?
\end{enumerate}

\exo[Une suite qui grandit vite]
Soit la suite $(u_{n})$ définie par $u_{0}=0$ et, pour tout $n\in\N$,
$u_{n+1}=u_{n}^{2}+1$.
\begin{enumerate}
  \item Calculer $u_{1}$, $u_{2}$, $u_{3}$, $u_{4}$ et $u_{5}$.
  \item Pour $n$ allant de $0$ à $5$, comparer $u_{n}$ et $2^{n}$.
  \item Démontrer par récurrence que, pour tout entier $n\geq 4$, on a
        $u_{n}\geq 2^{n}$ (on pourra remarquer que $4^{n}=2^{n}\times 2^{n}$).
\end{enumerate}

\partiefiche{Récurrence et terme général d'une suite}

\exo[Démontrer l'expression de $u_{n}$]
Pour chaque suite, calculer les deux termes qui suivent le premier, puis démontrer
par récurrence l'expression du terme général.
\begin{enumerate}[label=\alph*)]
  \item $(u_{n})$ définie par $u_{0}=5$ et, pour tout $n\in\N$, $u_{n+1}=2u_{n}-3$ :
        pour tout $n\in\N$, on a $u_{n}=2^{n+1}+3$.
  \item $(v_{n})$ définie par $v_{0}=-2$ et, pour tout $n\in\N$,
        $v_{n+1}=v_{n}+2n+4$ : pour tout $n\in\N$, on a $v_{n}=n^{2}+3n-2$.
  \item $(w_{n})$ définie par $w_{0}=80$ et, pour tout $n\in\N$,
        $w_{n+1}=1{,}2\,w_{n}-4$ : pour tout $n\in\N$, on a
        $w_{n}=60\times 1{,}2^{n}+20$.
\end{enumerate}

\exo[Des vélos en libre-service]
Une ville lance un service de vélos en libre-service en janvier 2026. Pour tout
$n\in\N$, on note $u_{n}$ le nombre d'abonnés $n$ mois après janvier 2026 : $u_{0}=300$.
Chaque mois, $20\,\%$ des abonnés résilient et $100$ nouveaux s'abonnent, donc, pour
tout $n\in\N$, $u_{n+1}=0{,}8\,u_{n}+100$.
\begin{enumerate}
  \item Calculer $u_{1}$ et $u_{2}$.
  \item Démontrer par récurrence que, pour tout $n\in\N$, on a
        $u_{n}=500-200\times 0{,}8^{n}$.
  \item Combien le service comptera-t-il d'abonnés en janvier 2027 ?
  \item Justifier que le nombre d'abonnés restera toujours inférieur à $500$.
  \item Compléter la fonction Python ci-dessous pour qu'elle renvoie le premier rang
        $n$ tel que $u_{n}>480$.
\begin{python}
def seuil():
    n = 0
    u = 300
    while ...:
        n = n + 1
        u = ...
    return n
\end{python}
  \item À l'aide de la calculatrice, déterminer le mois où le service dépassera
        $480$ abonnés.
\end{enumerate}

\exo[Conjecturer, puis démontrer]
Soit la suite $(u_{n})$ définie par $u_{0}=0$ et, pour tout $n\in\N$,
$u_{n+1}=4u_{n}+3$.
\begin{enumerate}
  \item Recopier et compléter le tableau.
\begin{center}
\renewcommand{\arraystretch}{1.3}
\begin{tabular}{|c|c|c|c|c|c|}
\hline
$n$ & $0$ & $1$ & $2$ & $3$ & $4$ \\ \hline
$u_{n}$ & $0$ & & & & \\ \hline
$u_{n}+1$ & $1$ & & & & \\ \hline
\end{tabular}
\end{center}
  \item Conjecturer l'expression de $u_{n}+1$, puis de $u_{n}$, en fonction de $n$.
  \item Démontrer cette conjecture par récurrence.
\end{enumerate}

\exo[Une suite homographique]
Soit la suite $(u_{n})$ définie par $u_{0}=\dfrac{1}{2}$ et, pour tout $n\in\N$,
$u_{n+1}=\dfrac{3u_{n}}{1+2u_{n}}$. On admet que $u_{n}>0$ pour tout $n\in\N$.
\begin{enumerate}
  \item Calculer $u_{1}$ et $u_{2}$ sous forme de fractions irréductibles.
  \item Conjecturer l'expression de $u_{n}$ en fonction de $n$ (on pourra comparer
        les numérateurs et les dénominateurs à des puissances de $3$).
  \item Démontrer cette conjecture par récurrence.
  \item Soit la suite $(v_{n})$ définie pour tout $n\in\N$ par
        $v_{n}=\dfrac{1}{u_{n}}-1$. Démontrer que $(v_{n})$ est géométrique, et
        retrouver l'expression de $u_{n}$.
\end{enumerate}

\exo[Expression et sens de variation]
Soit la suite $(u_{n})$ définie par $u_{0}=2$ et, pour tout $n\in\N$,
$u_{n+1}=2u_{n}-n$.
\begin{enumerate}
  \item Calculer $u_{1}$, $u_{2}$ et $u_{3}$.
  \item Démontrer par récurrence que, pour tout $n\in\N$, on a $u_{n}=2^{n}+n+1$.
  \item En déduire le sens de variation de la suite $(u_{n})$.
  \item Calculer $u_{20}$.
\end{enumerate}

\partiefiche{Récurrence et sens de variation}

\exo[Démontrer qu'une suite est monotone]
Pour chaque suite, démontrer par récurrence le sens de variation annoncé.
\begin{enumerate}[label=\alph*)]
  \item $(a_{n})$ définie par $a_{0}=1$ et, pour tout $n\in\N$,
        $a_{n+1}=0{,}5\,a_{n}+3$ : $(a_{n})$ est croissante.
  \item $(b_{n})$ définie par $b_{0}=10$ et, pour tout $n\in\N$,
        $b_{n+1}=0{,}5\,b_{n}+3$ : $(b_{n})$ est décroissante.
  \item $(c_{n})$ définie par $c_{0}=1$ et, pour tout $n\in\N$, $c_{n+1}=3c_{n}-1$ :
        $(c_{n})$ est croissante.
\end{enumerate}

\exo[Avec une racine carrée]
Soit la suite $(u_{n})$ définie par $u_{0}=1$ et, pour tout $n\in\N$,
$u_{n+1}=\sqrt{u_{n}+12}$. On admet que $u_{n}\geq 0$ pour tout $n\in\N$.
\begin{enumerate}
  \item Démontrer par récurrence que la suite $(u_{n})$ est croissante (on rappelle que
        la fonction racine carrée est croissante sur $\intfo{0}{+\infty}$).
  \item On remplace $u_{0}=1$ par $u_{0}=5$. Étudier le sens de variation de la suite
        obtenue.
  \item Que se passe-t-il si $u_{0}=4$ ?
\end{enumerate}

\exo[Une fonction et trois suites]
Soit $f$ la fonction définie sur $\R$ par $f(x)=x^{3}-6x^{2}+12x-6$.
\begin{enumerate}
  \item Calculer $f'(x)$, puis étudier le sens de variation de $f$ sur $\R$ (on pourra
        factoriser $f'(x)$).
  \item Soit la suite $(u_{n})$ définie par $u_{0}=1{,}5$ et, pour tout $n\in\N$,
        $u_{n+1}=f(u_{n})$. Démontrer par récurrence que, pour tout $n\in\N$, on a
        $u_{n}\leq u_{n+1}$. Qu'en déduit-on pour la suite $(u_{n})$ ?
  \item Étudier de même le sens de variation de la suite $(v_{n})$ définie par
        $v_{0}=2{,}5$ et, pour tout $n\in\N$, $v_{n+1}=f(v_{n})$.
  \item Soit la suite $(w_{n})$ définie pour tout $n\in\N$ par $w_{n}=f(n)$. Quel est
        son sens de variation ?
\end{enumerate}

\exo[Deux méthodes pour un même résultat]
Soit la suite $(u_{n})$ définie par $u_{0}=6$ et, pour tout $n\in\N$,
$u_{n+1}=3u_{n}-10$.
\begin{enumerate}
  \item Démontrer par récurrence que, pour tout $n\in\N$, on a $u_{n}\geq 5$.
  \item Exprimer $u_{n+1}-u_{n}$ en fonction de $u_{n}$, puis en déduire le sens de
        variation de $(u_{n})$.
  \item Retrouver ce sens de variation par une récurrence.
  \item Démontrer par récurrence que, pour tout $n\in\N$, on a $u_{n}=3^{n}+5$.
\end{enumerate}

\partiefiche{Récurrence et encadrement d'une suite}

\exo[Encadrer une suite par récurrence]
Pour chaque suite, démontrer l'encadrement par récurrence.
\begin{enumerate}[label=\alph*)]
  \item $(a_{n})$ définie par $a_{0}=1$ et, pour tout $n\in\N$,
        $a_{n+1}=0{,}5\,a_{n}+1$ : pour tout $n\in\N$, on a $1\leq a_{n}\leq 2$.
  \item $(b_{n})$ définie par $b_{0}=1$ et, pour tout $n\in\N$,
        $b_{n+1}=\sqrt{3b_{n}+4}$ : pour tout $n\in\N$, on a $1\leq b_{n}\leq 4$
        (la fonction racine carrée est croissante sur $\intfo{0}{+\infty}$).
  \item $(c_{n})$ définie par $c_{0}=3$ et, pour tout $n\in\N$,
        $c_{n+1}=\dfrac{4}{c_{n}+1}$ : pour tout $n\in\N$, on a $1\leq c_{n}\leq 3$
        (la fonction inverse est décroissante sur $\intoo{0}{+\infty}$).
\end{enumerate}

\exo[Encadrer une suite]
\textbf{Partie A.} Soit $f$ la fonction définie sur $\intfo{0}{+\infty}$ par
$f(x)=0{,}25x^{2}+1$ et la suite $(u_{n})$ définie par $u_{0}=0$ et, pour tout
$n\in\N$, $u_{n+1}=f(u_{n})$.
\begin{enumerate}
  \item Étudier le sens de variation de $f$ sur $\intfo{0}{+\infty}$.
  \item Démontrer par récurrence que, pour tout $n\in\N$, on a
        $0\leq u_{n}\leq u_{n+1}\leq 2$.
  \item En déduire le sens de variation de $(u_{n})$ et un encadrement de ses termes.
\end{enumerate}
\textbf{Partie B.} Soit la suite $(w_{n})$ définie par $w_{0}=1$ et, pour tout
$n\in\N$, $w_{n+1}=\sqrt{5w_{n}}$.
\begin{enumerate}[resume]
  \item Démontrer par récurrence que, pour tout $n\in\N$, on a
        $1\leq w_{n}\leq w_{n+1}\leq 5$.
  \item En déduire le sens de variation de $(w_{n})$.
\end{enumerate}

\exo[Une démonstration guidée]
Soit la suite $(v_{n})$ définie par $v_{0}=0$ et, pour tout $n\in\N$,
$v_{n+1}=\sqrt{v_{n}+6}$. Pour tout $n\in\N$, on note $P_{n}$ la propriété
«~$0\leq v_{n}\leq 3$~».
\begin{enumerate}
  \item Vérifier que $P_{0}$ est vraie.
  \item Soit $n\in\N$ tel que $P_{n}$ est vraie. Démontrer que $P_{n+1}$ est vraie.
  \item Que peut-on conclure des questions précédentes ?
  \item Démontrer de même que, pour tout $n\in\N$, on a
        $0\leq v_{n}\leq v_{n+1}\leq 3$, puis en déduire le sens de variation de
        $(v_{n})$.
\end{enumerate}

\exo[Des poissons dans un étang]
Au 1\textsuperscript{er} janvier 2026, un étang compte $7\,000$ poissons. Pour tout
$n\in\N$, on note $v_{n}$ le nombre de poissons, en millier, au 1\textsuperscript{er}
janvier $2026+n$. On modélise : $v_{0}=7$ et, pour tout $n\in\N$, $v_{n+1}=f(v_{n})$,
où $f$ est la fonction définie sur $\intff{0}{7{,}5}$ par $f(x)=-0{,}1x^{2}+1{,}5x$.
\begin{enumerate}
  \item Combien l'étang comptera-t-il de poissons au 1\textsuperscript{er} janvier 2027 ?
  \item Démontrer que la fonction $f$ est croissante sur $\intff{0}{7{,}5}$.
  \item Démontrer par récurrence que, pour tout $n\in\N$, on a
        $5\leq v_{n+1}\leq v_{n}\leq 7$.
  \item Interpréter ce résultat dans le contexte de l'exercice.
  \item À l'aide de la calculatrice, déterminer l'année à partir de laquelle l'étang
        comptera moins de $5\,100$ poissons.
\end{enumerate}

\exo[Encadrement et suite auxiliaire]
Soit $f$ la fonction définie sur $\intfo{0}{+\infty}$ par $f(x)=\dfrac{4x+3}{x+2}$ et la
suite $(u_{n})$ définie par $u_{0}=1$ et, pour tout $n\in\N$, $u_{n+1}=f(u_{n})$.
\begin{enumerate}
  \item Calculer $f'(x)$, puis étudier le sens de variation de $f$ sur
        $\intfo{0}{+\infty}$.
  \item Démontrer par récurrence que, pour tout $n\in\N$, on a
        $1\leq u_{n}\leq u_{n+1}\leq 3$.
  \item En déduire le sens de variation de $(u_{n})$.
  \item Soit la suite $(v_{n})$ définie pour tout $n\in\N$ par
        $v_{n}=\dfrac{u_{n}-3}{u_{n}+1}$. Démontrer que $(v_{n})$ est géométrique de
        raison $\frac{1}{5}$.
  \item En déduire l'expression de $u_{n}$ en fonction de $n$.
\end{enumerate}

\partiefiche{Sommes et récurrence}

\exo[Écrire une somme sans le symbole $\Sigma$]
Écrire chaque somme sans le symbole $\Sigma$, puis calculer les sommes a) à d).
\begin{center}
a)~$\displaystyle\sum_{k=1}^{6}k$ \qquad b)~$\displaystyle\sum_{k=0}^{4}k^{2}$ \qquad
c)~$\displaystyle\sum_{k=0}^{3}(3k+2)$ \\[4pt]
d)~$\displaystyle\sum_{k=2}^{5}\frac{1}{k}$ \qquad e)~$\displaystyle\sum_{k=0}^{n}2^{k}$
\end{center}

\exo[Écrire une somme avec le symbole $\Sigma$]
Écrire chaque somme avec le symbole $\Sigma$.
\begin{enumerate}[label=\alph*)]
  \item $1+2+3+\dots+50$
  \item $2+4+6+\dots+100$
  \item $1+3+5+\dots+(2n+1)$
  \item $1+q+q^{2}+\dots+q^{n}$
  \item $\dfrac{1}{1\times 2}+\dfrac{1}{2\times 3}+\dots+\dfrac{1}{n(n+1)}$
  \item $u_{3}+u_{4}+\dots+u_{n}$, pour un entier $n\geq 3$
\end{enumerate}

\exo[Démontrer une formule de somme]
Pour chaque formule, écrire la somme pour $n=3$ sans le symbole $\Sigma$ et vérifier
la formule, puis la démontrer par récurrence pour tout $n\in\Ns$.
\begin{enumerate}[label=\alph*)]
  \item $\displaystyle\sum_{k=1}^{n}(6k-2)=3n^{2}+n$
  \item $\displaystyle\sum_{k=1}^{n}(2k-1)=n^{2}$
  \item $\displaystyle\sum_{k=1}^{n}k^{2}=\frac{n(n+1)(2n+1)}{6}$
\end{enumerate}

\exo[La somme des cubes]
Pour tout $n\in\Ns$, on note $S_{n}=\displaystyle\sum_{k=1}^{n}k^{3}$.
\begin{enumerate}
  \item Calculer $S_{1}$, $S_{2}$, $S_{3}$ et $S_{4}$, puis les comparer aux carrés
        de $1$, $1+2$, $1+2+3$ et $1+2+3+4$.
  \item Rappeler l'expression de $1+2+\dots+n$ en fonction de $n$, et conjecturer
        celle de $S_{n}$.
  \item Démontrer cette conjecture par récurrence.
\end{enumerate}

\exo[Une somme de fractions]
Pour tout $n\in\Ns$, on note $T_{n}=\displaystyle\sum_{k=1}^{n}\frac{1}{k(k+1)}$.
\begin{enumerate}
  \item Calculer $T_{1}$, $T_{2}$ et $T_{3}$ sous forme de fractions irréductibles.
  \item Conjecturer l'expression de $T_{n}$ en fonction de $n$.
  \item Démontrer cette conjecture par récurrence.
  \item Vérifier que, pour tout $k\in\Ns$, on a
        $\dfrac{1}{k}-\dfrac{1}{k+1}=\dfrac{1}{k(k+1)}$. En écrivant $T_{n}$ terme à
        terme, retrouver le résultat.
\end{enumerate}

\exo[Somme des termes d'une suite géométrique]
\begin{enumerate}
  \item Démontrer par récurrence que, pour tout $n\in\N$, on a
        $\displaystyle\sum_{k=0}^{n}2^{k}=2^{n+1}-1$.
  \item Démontrer par récurrence que, pour tout $n\in\N$, on a
        $\displaystyle\sum_{k=0}^{n}3^{k}=\frac{3^{n+1}-1}{2}$.
  \item Soit $q$ un nombre réel différent de $1$. Démontrer par récurrence que, pour
        tout $n\in\N$, on a $\displaystyle\sum_{k=0}^{n}q^{k}=\frac{q^{n+1}-1}{q-1}$.
\end{enumerate}

\exo[Un tournoi de tennis de table]
Dans un tournoi de tennis de table, chaque joueur rencontre une fois chacun des
autres. Pour tout entier $n\geq 2$, on note $u_{n}$ le nombre de matchs d'un tournoi
à $n$ joueurs.
\begin{enumerate}
  \item Déterminer $u_{2}$, $u_{3}$ et $u_{4}$.
  \item Un nouveau joueur rejoint un tournoi à $n$ joueurs. Justifier que
        $u_{n+1}=u_{n}+n$.
  \item Démontrer par récurrence que, pour tout entier $n\geq 2$, on a
        $u_{n}=\dfrac{n(n-1)}{2}$.
  \item Écrire $u_{n}$ avec le symbole $\Sigma$.
  \item Combien de matchs compte un tournoi à $20$ joueurs ?
\end{enumerate}

\exo[Sommer les termes d'une suite]
Soit la suite $(u_{n})$ définie pour tout $n\in\N$ par $u_{n}=3-2\times 0{,}5^{n}$. Pour
tout $n\in\N$, on note $S_{n}=\displaystyle\sum_{k=0}^{n}u_{k}$.
\begin{enumerate}
  \item Calculer $u_{0}$, $u_{1}$ et $u_{2}$, puis $S_{0}$, $S_{1}$ et $S_{2}$.
  \item Compléter la fonction Python ci-dessous pour qu'elle renvoie $S_{n}$.
\begin{python}
def somme(n):
    s = 0
    for k in range(...):
        s = ...
    return s
\end{python}
  \item Démontrer par récurrence que, pour tout $n\in\N$, on a
        $S_{n}=3n-1+2\times 0{,}5^{n}$.
\end{enumerate}

\partiefichesansnum{Exercices type bac}

\exo[D'après bac 2025, Métropole]
Une équipe de biologistes étudie l'évolution de la superficie recouverte par une algue
marine appelée posidonie, sur le fond de la baie de l'Alycastre, près de l'île de
Porquerolles. La zone étudiée est d'une superficie totale de $20$ hectares (ha), et au
premier juillet 2024, la posidonie recouvrait $1$~ha de cette zone.

Pour tout entier naturel $n$, on note $u_{n}$ la superficie de la zone, en hectare,
recouverte par la posidonie au premier juillet de l'année $2024+n$. Ainsi, $u_{0}=1$.
Une étude conduite sur cette superficie a permis d'établir que pour tout entier naturel
$n$ : $u_{n+1}=-0{,}02\,u_{n}^{2}+1{,}3\,u_{n}$.
\begin{enumerate}
  \item Calculer la superficie que devrait recouvrir la posidonie au premier juillet
        2025 d'après ce modèle.
  \item On note $h$ la fonction définie sur $\intff{0}{20}$ par
        $h(x)=-0{,}02x^{2}+1{,}3x$. On admet que $h$ est croissante sur
        $\intff{0}{20}$. Démontrer que pour tout entier naturel $n$,
        $1\leq u_{n}\leq u_{n+1}\leq 20$.
  \item Les biologistes souhaitent savoir au bout de combien de temps la surface
        recouverte par la posidonie dépassera les $14$ hectares.
  \begin{enumerate}[label=\alph*)]
    \item Recopier et compléter l'algorithme suivant pour qu'en fin d'exécution, il
          affiche la réponse à la question des biologistes.
\begin{python}
def seuil():
    n = 0
    u = 1
    while ...... :
        n = ......
        u = ......
    return n
\end{python}
    \item À l'aide de la calculatrice, répondre à la question des biologistes.
  \end{enumerate}
\end{enumerate}

\exo[D'après bac 2025, Polynésie]
Une donnée binaire est une donnée qui ne peut prendre que deux valeurs : $0$ ou $1$.
Une donnée de ce type est transmise successivement d'une machine à une autre. Chaque
machine transmet la donnée reçue soit de manière fidèle, c'est-à-dire en transmettant
l'information telle qu'elle l'a reçue ($1$ devient $1$ et $0$ devient $0$), soit de façon
contraire ($1$ devient $0$ et $0$ devient $1$). La transmission est fidèle dans
$90\,\%$ des cas, et donc contraire dans $10\,\%$ des cas. Dans tout l'exercice, la
première machine reçoit toujours la valeur $1$.

Pour tout entier naturel $n\geq 1$, on note :
\begin{itemize}
  \item $V_{n}$ l'événement : «~la $n$-ième machine détient la valeur~$1$~» ;
  \item $\overline{V_{n}}$ l'événement : «~la $n$-ième machine détient la valeur~$0$~».
\end{itemize}
\begin{enumerate}
  \item
  \begin{enumerate}[label=\alph*)]
    \item Recopier et compléter l'arbre de probabilité ci-dessous.
\begin{center}
\begin{tikzpicture}[grow=right,level distance=1.9cm,font=\small,
    level 1/.style={sibling distance=1.6cm},level 2/.style={sibling distance=0.8cm}]
  \node {$V_{1}$}
    child {node {$\overline{V_{2}}$}
      child {node {$\overline{V_{3}}$} edge from parent node[below,font=\scriptsize]{$\dots$}}
      child {node {$V_{3}$} edge from parent node[above,font=\scriptsize]{$\dots$}}
      edge from parent node[below,font=\scriptsize]{$\dots$}}
    child {node {$V_{2}$}
      child {node {$\overline{V_{3}}$} edge from parent node[below,font=\scriptsize]{$\dots$}}
      child {node {$V_{3}$} edge from parent node[above,font=\scriptsize]{$\dots$}}
      edge from parent node[above,font=\scriptsize]{$\dots$}};
\end{tikzpicture}
\end{center}
    \item Démontrer que $P(V_{3})=0{,}82$ et interpréter ce résultat dans le contexte
          de l'exercice.
    \item Sachant que la troisième machine a reçu la valeur $1$, calculer la
          probabilité que la deuxième machine ait aussi reçu la valeur $1$.
  \end{enumerate}
  \item Pour tout entier naturel $n\geq 1$, on note $p_{n}=P(V_{n})$. La première
        machine a reçu la valeur $1$, on a donc $p_{1}=1$.
  \begin{enumerate}[label=\alph*)]
    \item Démontrer que pour tout entier naturel $n\geq 1$ :
          $p_{n+1}=0{,}8\,p_{n}+0{,}1$.
    \item Démontrer par récurrence que pour tout entier naturel $n\geq 1$,
          $p_{n}=0{,}5\times 0{,}8^{n-1}+0{,}5$.
  \end{enumerate}
\end{enumerate}

\exo[D'après bac 2026, Métropole]
\textbf{Partie A.} Soit $f$ la fonction définie sur $\intoo{-3}{+\infty}$ par
$f(x)=\dfrac{2x}{x+3}$. On admet que $f$ est dérivable sur $\intoo{-3}{+\infty}$.
\begin{enumerate}
  \item Vérifier, en détaillant les calculs, que pour tout nombre réel $x$ de
        $\intoo{-3}{+\infty}$, $f'(x)=\dfrac{6}{(x+3)^{2}}$. En déduire que $f$ est
        croissante sur $\intoo{-3}{+\infty}$.
\end{enumerate}
Soit $(u_{n})$ la suite définie par $u_{0}=1$ et pour tout entier naturel $n$ par
$u_{n+1}=f(u_{n})$.
\begin{enumerate}[resume]
  \item Démontrer, par récurrence, que pour tout entier $n$, on a :
        $0\leq u_{n+1}\leq u_{n}\leq 1$.
\end{enumerate}
\textbf{Partie B.} On admet que pour tout entier naturel $n$, $u_{n}\neq 0$. On note
$(v_{n})$ et $(w_{n})$ les suites définies pour tout entier naturel $n$ par
$v_{n}=\dfrac{1}{u_{n}}$ et $w_{n}=v_{n}+1$.
\begin{enumerate}
  \item Vérifier que, pour tout entier naturel $n$,
        $v_{n+1}=\dfrac{3}{2}v_{n}+\dfrac{1}{2}$.
  \item Démontrer que $(w_{n})$ est une suite géométrique dont on précisera les
        paramètres.
  \item En déduire que, pour tout entier naturel $n$,
        $u_{n}=\dfrac{1}{2\left(\frac{3}{2}\right)^{n}-1}$.
  \item
  \begin{enumerate}[label=\alph*)]
    \item Recopier et compléter le script suivant de telle sorte que la fonction
          \texttt{seuil()} renvoie le plus petit entier $n$ tel que $u_{n}$ soit
          inférieur ou égal à $10^{-3}$.
\begin{python}
def seuil():
    n = 0
    u = 1
    while ... :
        n = n + 1
        u = ...
    return ...
\end{python}
    \item Quelle est la valeur renvoyée par la fonction \texttt{seuil()} ?
  \end{enumerate}
\end{enumerate}

\end{multicols}

\end{document}
